% ===================================================================
%  TABLEAU N° 8 — La Moisson. — Les aspects de la campagne. —
%  La Chasse. — La Vendange. — La Pêche.
%  Feuillets 51 a 55 du fac-simile = folios 49 a 53.
%  Correspondance : folio imprime = numero de feuillet - 2.
%
%  Les cles « %%K » sont les MEMES que dans texto/io/17-tabelo-08.tex.
%  Ce tableau combine deux « scenes » : la Premiere scene (La Moisson :
%  les aspects de la campagne, le pastre, la moisson propre, l'orage)
%  et la Deuxieme scene (La Chasse, la Vendange, la Biere et le Cidre,
%  la Peche). La numerotation des alineas et des renvois (1), (2)...
%  est propre a chaque scene et reprend a 1 a la Deuxieme scene — d'ou
%  la repetition des cles %%K t08-01, t08-02... pour la deuxieme
%  scene : c'est l'imprime lui-meme qui recommence, non le releve.
%
%  ATTENTION — LE FAC-SIMILE FRANCAIS EST UNE PRISE DE VUE, contraste
%  inegal d'une page a l'autre ; plusieurs renvois entre parentheses
%  ont ete verifies par recoupement avec les memes renvois du livret
%  ido (meme illustration, memes numeros de legende dans les deux
%  editions) apres agrandissement du fac-simile.
% ===================================================================

% --- Feuillet 51 (folio 49) -----------------------------------------
\begin{VUpage}[51]{49}
%%K t08-apar-1 apar
\VUsaut{5.0mm}
\VUtitre{15.0pt}{60}{TABLEAU N\textsuperscript{o} 8}
\VUsaut{1.4mm}
\VUcentre{13.2pt}{90}{Première scène.}
\VUsaut{2.2mm}
\VUtitre{11.4pt}{0}{La Moisson. --- Les aspects de la campagne.}
\VUsaut{6.4mm}

%%K t08-c1-01-1 p
\VUblancAlinea
1. --- Avec l'été, l'aspect de la campagne change\nl
encore une fois. La semence confiée au sillon a germé ;\nl
une petite plante verte est sortie de terre et n'a pas\nl
tardé à porter un épi. Le grain s'est formé, puis il a\nl
mûri et il ne reste plus qu'à moissonner.

%%K t08-c1-02-1 p
\VUblancAlinea
2. --- Les \VUgras{fleurs des champs} sont épanouies.\nl
\VUgras{Bluets} \textsuperscript{(1)}, \VUgras{coquelicots} \textsuperscript{(2)} et \VUgras{digitales} \textsuperscript{(3)} mettent\nl
dans la teinte uniforme des blés la note gaie de leurs\nl
fraîches \VUgras{corolles.}

%%K t08-c1-03-1 p
\VUblancAlinea
3. --- Les arbres fruitiers se sont garnis de feuilles,\nl
le fruit est apparu. Le petit \VUgras{cerisier} \textsuperscript{(4)} est couvert\nl
de belles \VUgras{cerises} \textsuperscript{(5)}, que les enfants cueillent et man\cc
gent avec délices.

%%K t08-c1-04-1 p
\VUblancAlinea
4. --- Le troupeau peut passer à présent toute la\nl
journée en plein air.

%%K t08-c1-04-2 p
\VUblancAlinea
Les \VUgras{agneaux} \textsuperscript{(6)} ont grandi et sont devenus de belles\nl
\VUgras{brebis} \textsuperscript{(7)} et de beaux \VUgras{moutons} \textsuperscript{(8)} couverts d'une\nl
épaisse toison. Ils broutent paisiblement l'\VUgras{herbe} du\nl
\VUgras{pâturage} \textsuperscript{(9)} parsemée de \VUgras{pâquerettes.}

%%K t08-c1-04-3 p
\VUblancAlinea
On va bientôt commencer à tondre ces animaux\nl
pour avoir leur \VUgras{laine}, avec laquelle on fabrique le\nl
drap de nos vêtements.

%%K t08-c1-05-1 p
\VUblancAlinea
5. --- Le \VUgras{pâtre} \textsuperscript{(10)} ne paraît guère s'occuper d'eux.\nl
Il ne s'inquiète que de l'orage qui passe à l'horizon, et\nl
le \VUgras{berger} \textsuperscript{(13)}, qui porte sa \VUgras{gourde} \textsuperscript{(14)} à ses lèvres,\nl
semble encore plus insouciant.

%%K t08-c1-06-1 p
\VUblancAlinea
6. --- La chaleur est accablante ; car le \VUgras{chien} \textsuperscript{(15)}\nl
vient lui aussi se rafraîchir ; il lappe l'eau fraîche\nl
du \VUgras{ruisseau} \textsuperscript{(16)} et effraie une pauvre petite gre\cc
nouille \textsuperscript{(17)} qui était tapie dans l'herbe du rivage.
\parplein
\end{VUpage}

% --- Feuillet 52 (folio 50) -----------------------------------------
\begin{VUpage}[52]{50}
%%K t08-c1-06-1 p suite
\VUcontinue
Elle saute dans l'eau claire où s'entr'ouvrent les\nl
\VUgras{nénuphars} \textsuperscript{(18)}.

%%K t08-c1-07-1 p
\VUblancAlinea
7. --- Une \VUgras{bergeronnette} \textsuperscript{(19)} se baigne près de la\nl
petite \VUgras{source} \textsuperscript{(20)} qui jaillit entre deux rochers et se\nl
cache sous les \VUgras{broussailles} \textsuperscript{(21)} et les \VUgras{buissons}\nl
\VUgras{d'aubépine} \textsuperscript{(22)}.

%%K t08-c1-08-1 p
\VUblancAlinea
8. --- Pour les enfants de la campagne, la moisson\nl
est une époque très amusante. Ils font de joyeuses cul\cc
butes \textsuperscript{(24)} sur les \VUgras{tas de paille} \textsuperscript{(23)}, ils aident les\nl
\VUgras{moissonneurs} \textsuperscript{(26)}, et, tandis que ceux-ci fauchent\nl
le \VUgras{champ de blé} \textsuperscript{(27)} avec leurs \VUgras{faux} \textsuperscript{(25)} bien aiguisées\nl
sur la \VUgras{pierre} \textsuperscript{(30)}, ils ramassent le blé et le lient en\nl
\VUgras{gerbes} \textsuperscript{(31)} que l'on réunit ensuite en \VUgras{javelles} \textsuperscript{(32)}.

%%K t08-c1-09-1 p
\VUblancAlinea
9. --- On permet quelquefois aux plus grands de\nl
couper avec une \VUgras{faucille} \textsuperscript{(12)} les \VUgras{tiges} \textsuperscript{(28)} dorées qui\nl
portent les lourds \VUgras{épis} \textsuperscript{(29)} remplis de grain, et les\nl
petits, tout en gardant le \VUgras{bétail} \textsuperscript{(33)} qui broute tran\cc
quillement dans la \VUgras{prairie} \textsuperscript{(34)}, à l'ombre du grand\nl
\VUgras{tilleul} \textsuperscript{(35)}, attrapent dans leur \VUgras{filet} \textsuperscript{(36)} les beaux\nl
papillons.

%%K t08-c1-09-2 p
\VUblancAlinea
Quelques-uns même grimpent sur la \VUgras{charrette} \textsuperscript{(37)}\nl
pour ranger les gerbes qu'on leur passe au bout d'une\nl
\VUgras{fourche} \textsuperscript{(38)}.

%%K t08-c1-10-1 p
\VUblancAlinea
10. --- Dans toute la \VUgras{plaine} \textsuperscript{(39)} où s'envolent les\nl
\VUgras{alouettes} \textsuperscript{(41)}, l'animation règne. Tout le monde est\nl
occupé à la moisson. Mais, tandis que les pauvres\nl
gens continuent à se servir des vieux instruments,\nl
\VUgras{faux,} \VUgras{faucilles} et \VUgras{fléaux} \textsuperscript{(40)}, les riches propriétaires\nl
font faucher leur blé ou leur \VUgras{seigle} \textsuperscript{(43)} avec une\nl
\VUgras{moissonneuse mécanique} \textsuperscript{(42)}. Puis, sans avoir\nl
besoin de porter les gerbes dans la grange, on forme\nl
sur place de grandes \VUgras{meules} \textsuperscript{(44)}.

%%K t08-c1-11-1 p
\VUblancAlinea
11. --- On amène alors une \VUgras{machine à battre} \textsuperscript{(47)},\nl
qu'une \VUgras{locomobile} \textsuperscript{(45)} met en mouvement à l'aide\nl
d'une \VUgras{courroie} \textsuperscript{(46)} de transmission, et le grain se\nl
trouve ainsi séparé de la paille sans quitter le champ\nl
où il avait été semé.
\end{VUpage}

% --- Feuillet 53 (folio 51) -----------------------------------------
\begin{VUpage}[53]{51}
%%K t08-c1-12-1 p
\VUblancAlinea
12. --- Fréquemment, de violents \VUgras{orages} \textsuperscript{(53)}\nl
éclatent à l'époque de la moisson. Tout d'abord les\nl
grondements du \VUgras{tonnerre} se font entendre au loin ;\nl
un violent \VUgras{vent d'ouest} \textsuperscript{(54)} s'élève et fait plier en\nl
gémissant les plus gros arbres du petit \VUgras{bois} \textsuperscript{(49)}. De\nl
gros \VUgras{nuages} \textsuperscript{(56)} noirs montent dans le ciel ; ils sont\nl
sillonnés par les zigzags aveuglants de l'\VUgras{éclair} \textsuperscript{(55)}.\nl
La \VUgras{pluie} et parfois la \VUgras{grêle} se mettent à tomber,\nl
écrasant la moisson prête à être fauchée.

%%K t08-c1-13-1 p
\VUblancAlinea
13. --- Souvent la foudre incendie quelque bâtiment.\nl
C'est ainsi que l'année dernière le toit du \VUgras{moulin à}\nl
\VUgras{vent} \textsuperscript{(50)} fut réduit en cendres et que deux de ses\nl
\VUgras{ailes} \textsuperscript{(51)} furent pulvérisées.

%%K t08-c1-14-1 p
\VUblancAlinea
14. --- Puis, au bout de quelques heures, le calme\nl
renaît. Un brillant \VUgras{arc-en-ciel} \textsuperscript{(52)} se montre à l'ho\cc
rizon et la nature reprend sa vie un instant interrom\cc
pue. Les paysans, qui se sont réfugiés dans les caba\cc
nes \textsuperscript{(48)}, se remettent au travail et essaient courageu\cc
sement de réparer le dégât qu'ils viennent de subir.

%%K t08-tit-1 sub
\VUsaut{5.4mm}
\VUcentre{13.2pt}{90}{Deuxième scène.}
\VUsaut{2.2mm}

%%K t08-tit-2 sub
\VUtitre{11.4pt}{0}{La Chasse. --- La Vendange. --- La Pêche.}
\VUsaut{4.0mm}

%%K t08-c2-01-1 p
\VUblancAlinea
1. --- Vers la fin d'\VUgras{août} ou vers le milieu de sep\cc
tembre, suivant les régions, a lieu l'ouverture de la\nl
chasse. A peine les premiers rayons du soleil ont-ils\nl
fait sortir les petits \VUgras{lézards} \textsuperscript{(2)} de leur trou, que le\nl
\VUgras{chasseur} \textsuperscript{(5)} se met en route avec ses \VUgras{chiens} \textsuperscript{(4)}.

%%K t08-c2-02-1 p
\VUblancAlinea
2. --- Il a endossé son solide \VUgras{costume} \textsuperscript{(9)} de chasse\nl
en gros drap, chaussé ses \VUgras{guêtres} de cuir, avec les\cc
quelles il pourra marcher dans l'herbe mouillée, où se\nl
cachent les \VUgras{crapauds} \textsuperscript{(1)}, il a pris son \VUgras{fusil} \textsuperscript{(6)} et\nl
son \VUgras{carnier} \textsuperscript{(10)}, et le voilà parti.

%%K t08-c2-03-1 p
\VUblancAlinea
3. --- L'ardeur de la poursuite l'amène au milieu d'un\nl
vignoble où la vendange bat son plein. Les enfants du\nl
vigneron, qui gardent ordinairement les chèvres sur\nl
la route, les laissent aujourd'hui brouter à l'aventure,
\parplein
\end{VUpage}

% --- Feuillet 54 (folio 52) -----------------------------------------
\begin{VUpage}[54]{52}
%%K t08-c2-03-1 p suite
\VUcontinue
et après avoir attaché à un pieu un vieux \VUgras{bouc} \textsuperscript{(3)} qui\nl
ne manquerait pas de profiter de sa liberté pour\nl
s'échapper, ils se donnent, eux aussi, leur part de\nl
plaisir. L'un prend une grappe de raisin dans un\nl
\VUgras{panier à vendange} \textsuperscript{(12)} et s'amuse à faire couler le\nl
\VUgras{jus} \textsuperscript{(14)} dans un \VUgras{gobelet} \textsuperscript{(13)} pour faire du vin doux.\nl
L'autre se coiffe d'une \VUgras{couronne de feuillage} \textsuperscript{(15)}\nl
qu'il vient de tresser.

%%K t08-c2-03-2 p
\VUblancAlinea
A côté d'eux, les \VUgras{moineaux} \textsuperscript{(16)} effrontés viennent\nl
jusque devant le pressoir picorer les grains de raisin.

%%K t08-c2-04-1 p
\VUblancAlinea
4. --- Pendant que les enfants s'amusent les hommes\nl
travaillent. Les \VUgras{vendangeurs} \textsuperscript{(18)} apportent le raisin\nl
au cellier dans des \VUgras{hottes} \textsuperscript{(17)} ; un \VUgras{tonnelier} \textsuperscript{(19)}\nl
répare les \VUgras{barriques} \textsuperscript{(20)}, vérifie si les \VUgras{douves} \textsuperscript{(22)} et\nl
les \VUgras{fonds} \textsuperscript{(21)} sont en bon état et remplace les cer\cc
cles \textsuperscript{(23)} cassés. C'est lui qui a fait les grandes \VUgras{cuves} \textsuperscript{(25)}\nl
dans lesquelles on verse le \VUgras{raisin} \textsuperscript{(26)} pour le faire\nl
fermenter.

%%K t08-c2-05-1 p
\VUblancAlinea
5. --- Lorsque la fermentation a eu lieu, on met le\nl
raisin sur le \VUgras{pressoir} \textsuperscript{(28)}, et, en serrant progressive\cc
ment la \VUgras{vis} \textsuperscript{(29)}, on exprime le jus, qui coule dans un\nl
\VUgras{baquet} \textsuperscript{(32)}, et l'on obtient le \VUgras{vin} \textsuperscript{(31)}.

%%K t08-c2-06-1 p
\VUblancAlinea
6. --- Sur le mur du \VUgras{cellier} \textsuperscript{(27)} grimpe une branche\nl
de \VUgras{houblon} \textsuperscript{(30)}. Ici ce n'est qu'une plante d'ornement,\nl
mais dans certains pays, où la vigne est très rare, le\nl
houblon est cultivé pour la fabrication de la \VUgras{bière.}

%%K t08-c2-07-1 p
\VUblancAlinea
7. --- Le petit \VUgras{pommier} \textsuperscript{(33)} est chargé de pommes\nl
qui feront, lorsqu'elles seront mûres, un \VUgras{cidre} excel\cc
lent. Dans leur \VUgras{cage} \textsuperscript{(34)}, le \VUgras{serin} \textsuperscript{(35)} et le \VUgras{chardon}\cc
\VUgras{neret} \textsuperscript{(36)} contemplent d'un œil d'envie, les moineaux\nl
qui s'ébattent en liberté.

%%K t08-c2-08-1 p
\VUblancAlinea
8. --- A perte de vue, sur le penchant des collines,\nl
s'alignent les \VUgras{échalas} \textsuperscript{(40)}, supportant les magnifiques\nl
\VUgras{ceps de vigne} \textsuperscript{(39)}, garnis de lourdes \VUgras{grappes.}

%%K t08-c2-08-2 p
\VUblancAlinea
Pendant toute l'année, le \VUgras{vigneron} \textsuperscript{(42)} a travaillé\nl
pour soigner son \VUgras{vignoble} \textsuperscript{(38)}, et maintenant il est\nl
récompensé de sa peine par une belle récolte. Les
\parplein
\end{VUpage}

% --- Feuillet 55 (folio 53) -----------------------------------------
\begin{VUpage}[55]{53}
%%K t08-c2-08-2 p suite
\VUcontinue
\VUgras{vendangeuses} \textsuperscript{(41)} emplissent les cuves placées sur\nl
la charrette traînée par des \VUgras{mules} \textsuperscript{(43)} et tout le monde\nl
est joyeux.

%%K t08-c2-09-1 p
\VUblancAlinea
9. --- Il en est de même des \VUgras{pêcheurs} \textsuperscript{(44)} à la ligne\nl
qui, dès le matin, sont venus s'installer sur le bord de\nl
l'eau avec leurs \VUgras{cannes à pêche} \textsuperscript{(45)}.

%%K t08-c2-09-2 p
\VUblancAlinea
Le \VUgras{poisson} \textsuperscript{(47)} mord à l'hameçon dès qu'ils jet\cc
tent leur \VUgras{ligne} \textsuperscript{(46)} dans l'eau, et à peine ont-ils le\nl
temps de renouveler l'\VUgras{amorce} qu'une nouvelle vic\cc
time frétille au bout du fil.

%%K t08-c2-09-3 p
\VUblancAlinea
Pourtant le \VUgras{pêcheur au filet} \textsuperscript{(48)}, qui, monté dans\nl
sa \VUgras{barque} \textsuperscript{(49)}, jette l'épervier au milieu de la\nl
\VUgras{rivière} \textsuperscript{(50)}, en prend bien davantage.

%%K t08-c2-10-1 p
\VUblancAlinea
10. --- L'automne est la saison des bonnes promena\cc
des : promenade à pied, à \VUgras{cheval} \textsuperscript{(52)}, à \VUgras{bicyclette} \textsuperscript{(53)},\nl
en \VUgras{automobile} \textsuperscript{(54)}. Les \VUgras{routes} \textsuperscript{(51)} sont propres et l'on\nl
profite des derniers beaux jours, car voilà déjà les\nl
\VUgras{hirondelles} \textsuperscript{(56)} qui se rassemblent sur les toits pour\nl
s'envoler à- tire-d'aile vers des pays plus chauds. Les\nl
feuilles du vieux \VUgras{noyer} \textsuperscript{(57)} jaunissent, les \VUgras{noix} tom\cc
bent et les grands \VUgras{arbres} groupés en \VUgras{bouquet} \textsuperscript{(58)}\nl
sur la \VUgras{hauteur} \textsuperscript{(59)} vont bientôt se dépouiller.

%%K t08-c2-11-1 p
\VUblancAlinea
11. --- Le \VUgras{soleil} \textsuperscript{(60)} se lève plus tard, les jours\nl
diminuent, un petit froid assez vif se fait sentir le\nl
\VUgras{matin} et voilà déjà des vols de \VUgras{bécasses} \textsuperscript{(61)} qui\nl
quittent nos climats inhospitaliers.

\VUsaut{6.0mm}
\VUfilet{22mm}
\end{VUpage}
